% Research-only fragment, NOT inserted into main.tex.
% Requires amsmath, amssymb, amsthm; environments theorem, proposition, corollary, remark, proof.
\section{Weighted first-excess transport beyond the coupling horizon}
\label{sec:first-excess-transport}

\paragraph{Context and notation.}
Fix an odd $c\ge-1$ with $3\nmid c$, and consider the accelerated map
\[
 T_{3,c}(n)=\begin{cases}n/2,&n\equiv0\pmod2,\\
 (3n+c)/2,&n\equiv1\pmod2.
 \end{cases}
\]
We use the precise first-passage observable and logarithmic averaging
of the accompanying manuscript.  Set $\beta=1/200$ and, for $X>1$,
\[
 a_X=X^{1-2\beta},\quad
 W_X=(X^{1-\beta},X]\cap\mathbb N,\quad
 H_X=X(\log X)^{\varepsilon_0/2},\quad
 L_X=H_X/X.
\]
For each odd $m$, write $b_1(m),\ldots,b_{s(m)}(m)$ for the
\emph{chronologically ordered} elements of $W_X$ encountered in the
Syracuse blocks up to and including the first passage below $a_X$.
As in the manuscript, the entire list is declared empty if that
first passage never occurs.  The first term $m$ itself is not included;
thus
\[
  \Phi_X(m)=\sum_{i=1}^{s(m)}b_i(m),\qquad
  Y_X(m)=\Phi_X(m)/X.
\]
Write $m=w^\flat$ for the odd part of a positive integer $w$, and let
$\mathbb E_{\log}$ denote the limiting logarithmic mean over $w$.
Throughout, we assume the manuscript's bounded-observable reduction
lemma and the proven-in-the-manuscript tail estimate
\begin{equation}\label{eq:fe-tightness}
 \mu_X(L):=\mathbb E_{\log}{\bf1}_{\{Y_X>L\}}
 \le C e^{-L/C}+C(\log X)^{-\varepsilon_0},\qquad L\ge0,
\end{equation}
with $C$ independent of $X,L$.  No unproved mixing or independence
assumption after the coupling horizon is introduced below.

\paragraph{The first-excess law.}
Fix $L\ge0$. If $Y_X(m)>L$, let
\[
 \kappa_{X,L}(m):=\min\Big\{k:\sum_{i=1}^{k}\frac{b_i(m)}X>L\Big\},
 \qquad J_{X,L}(m):=b_{\kappa_{X,L}(m)}(m),
\]
and otherwise set $J_{X,L}(m)=\dagger$.
For $y\in W_X$ define
\begin{equation}\label{eq:fe-entrance}
 \eta_{X,L}(y):=\mathbb E_{\log}{\bf1}_{\{J_{X,L}=y\}},
 \qquad p_X(L):=\sum_{y\in W_X}\eta_{X,L}(y)=\mu_X(L).
\end{equation}
All these limits exist under the stated reduction lemma: for fixed
$X,L,y$ the relevant indicator is bounded and unchanged when an
initial segment above $X$ is removed by an admissible high first
passage.  The state space $W_X$ is finite.

If $y\in W_X$ has a finite forward first passage below $a_X$, define
its \emph{post-state weighted occupation}
\begin{equation}\label{eq:fe-future-mass}
 r_X(y):=\frac1X\sum_{\substack{t\ge1:\,t<\tau_{\le a_X}(y)\\
                        T_{3,c}^{t}(y)\in W_X}}T_{3,c}^{t}(y),
\end{equation}
where $\tau_{\le a_X}(y)$ denotes the first accelerated-map time
at which the value is at most $a_X$.  Set $r_X(y)=0$ for states
without such a finite first passage.  On the support of
$\eta_{X,L}$ the first passage is finite, so this convention is
irrelevant.  Although the manuscript uses Syracuse blocks, counting
future members of $W_X$ until the accelerated map first reaches
$[1,a_X]$ gives the same remaining list: after an even term has
fallen below $a_X$, all subsequent members of the current final
block are still below $a_X<X^{1-\beta}$.

\begin{theorem}[Exact weighted first-excess transport]
\label{thm:first-excess-identity}
Under the preceding definitions, the limit
\[
  T_X(L):=\mathbb E_{\log}\bigl[(\Phi_X-LX)_+\bigr]
\]
exists and
\begin{equation}\label{eq:fe-exact}
 \boxed{\quad
 \frac{T_X(L)}X=\sum_{y\in W_X}\eta_{X,L}(y)r_X(y)
                    +\vartheta_X(L),\qquad
 0\le\vartheta_X(L)\le p_X(L).
 \quad}
\end{equation}
In particular, the only term not controlled by the probability
$p_X(L)$ is the \emph{future weighted mass transported from the
first-excess entrance law}.

Let $\mathcal S_X\subseteq W_X$ be the states that reach $[1,a_X]$.
Define the substochastic deterministic next-visit kernel $Q_X$ on
$\mathcal S_X$: $Q_X(y,z)=1$ when $z$ is the next visit to $W_X$
after $y$ and before that passage, and $Q_X(y,z)=0$ otherwise.
For $v_X(z):=z/X$ one also has the \emph{finite Green expansion}
\begin{equation}\label{eq:fe-green}
 \sum_{y\in W_X}\eta_{X,L}(y)r_X(y)
 =\sum_{k\ge1}\sum_{z\in\mathcal S_X}
            v_X(z)\,(\eta_{X,L}Q_X^k)(z).
\end{equation}
The sum is finite for each fixed $X$.
\end{theorem}

\begin{proof}
For a successful stopped orbit with $Y_X>L$, let $k=\kappa_{X,L}$.
Then
\[
 Y_X-L
 =\underbrace{\Big(\sum_{i=1}^{k}b_i/X-L\Big)}_{e_{X,L}\in(0,1]}
  +\sum_{i=k+1}^{s(m)}b_i/X.
\]
The last sum is $r_X(J_{X,L}(m))$, because the accelerated-map
trajectory after the selected state is deterministic and its stopping
threshold is fixed.  On $Y_X\le L$ both terms are zero.  Thus the
identity holds \emph{pointwise}. The overshoot $e_{X,L}$ is bounded
by $1$, since every $b_i\le X$.  All involved observables have
logarithmic means: the state $J_{X,L}$ takes finitely many values,
and the future mass $r_X(y)$ is finite on its support.  Taking means
proves \eqref{eq:fe-exact} and the bounds on $\vartheta_X(L)$.

A state in $\mathcal S_X$ cannot be visited twice before the stated
first passage, since a repetition would trap the deterministic orbit
in a periodic cycle disjoint from $[1,a_X]$.  As $\mathcal S_X$ is
finite, $Q_X$ is nilpotent.  Each subsequent $W_X$-visit from $y$
is one of the states represented by $Q_X^k(y,\cdot)$, proving
\eqref{eq:fe-green}.
\end{proof}

\paragraph{A scale-dependent allowance for later returns.}
Let
\[
 I_{X,j}=(X2^{-j-1},X2^{-j}]\cap W_X,\qquad j\ge0,
\]
which is nonempty only for finitely many $j$. For $y\in\mathcal S_X$,
let $N^+_{X,j}(y)$ denote the number of visits to $I_{X,j}$
\emph{strictly after} the state $y$ and before the first passage
below $a_X$.  For arbitrary integers $K_{X,j}\ge1$ set
\begin{equation}\label{eq:fe-excess-return}
 \mathcal E_X(L;K):=
  \sum_{y\in W_X}\eta_{X,L}(y)
       \sum_{j\ge0}2^{-j}\bigl(N^+_{X,j}(y)-K_{X,j}\bigr)_+.
\end{equation}

\begin{theorem}[Late-return localisation with a free weighted budget]
\label{thm:fe-free-budget}
For any $L\ge0$ and any such integer allowances,
\begin{equation}\label{eq:fe-sandwich}
 \boxed{\quad
 \frac12\mathcal E_X(L;K)
 \le \frac{T_X(L)}X
 \le p_X(L)\Bigl(1+\sum_{j\ge0}2^{-j}K_{X,j}\Bigr)
      +\mathcal E_X(L;K).
 \quad}
\end{equation}
Now fix $0<\sigma<\varepsilon_0$ and $0\le\alpha<1$, and take
\begin{equation}\label{eq:fe-allowances}
 L=L_X=(\log X)^{\varepsilon_0/2},\qquad
 K_{X,j}=\Big\lceil(\log X)^\sigma 2^{\alpha j}\Big\rceil.
\end{equation}
Then, with $T_X=T_X(L_X)$ as in the manuscript,
\begin{equation}\label{eq:fe-equivalence}
 \boxed{\quad
 T_X=o(X)\quad\Longleftrightarrow\quad
 \mathcal E_X(L_X;K)=o(1).
 \quad}
\end{equation}
In particular, the first $\asymp(\log X)^\sigma2^{\alpha j}$
\emph{post-excess visits in shell $j$} make a combined contribution
$o(X)$ to $T_X$, uniformly without any additional dynamical
hypothesis.
\end{theorem}

\begin{proof}
For $z\in I_{X,j}$,
$2^{-j-1}<z/X\le2^{-j}$.  Thus
\[
 \frac12\sum_{j\ge0}2^{-j}N^+_{X,j}(y)
 \le r_X(y)\le
 \sum_{j\ge0}2^{-j}N^+_{X,j}(y).
\]
Moreover $N\le K+(N-K)_+$ for $N,K\ge0$.  Insert these inequalities
into \eqref{eq:fe-exact}, using $0\le\vartheta_X(L)\le p_X(L)$;
this proves \eqref{eq:fe-sandwich}.

By \eqref{eq:fe-tightness},
\[
 p_X(L_X)\le Ce^{-L_X/C}+C(\log X)^{-\varepsilon_0}
       =O((\log X)^{-\varepsilon_0}).
\]
The chosen allowances satisfy
\[
 \sum_{j\ge0}2^{-j}K_{X,j}
 \le \frac{(\log X)^\sigma}{1-2^{\alpha-1}}+2
 =O_\alpha((\log X)^\sigma),
\]
so the first term on the right of \eqref{eq:fe-sandwich} is
$O((\log X)^{\sigma-\varepsilon_0})=o(1)$.
The two inequalities of \eqref{eq:fe-sandwich} now give the exact
equivalence \eqref{eq:fe-equivalence}.
\end{proof}

\begin{proposition}[An exact induced-shell transport formula]
\label{prop:fe-shell-transport}
For each nonempty $I_{X,j}$ let $E_{X,j}(y,z)$ be the
(first-future-entry) kernel from $\mathcal S_X$ to
$I_{X,j}\cap\mathcal S_X$: it is $1$ when $z$ is the first visit to
$I_{X,j}$ strictly after $y$ and before the terminal passage, and
$0$ otherwise.  Let $Q_{X,j}$ be the first-\emph{return} kernel on
$I_{X,j}\cap\mathcal S_X$, defined analogously.  Both are
substochastic deterministic kernels, and $Q_{X,j}$ is nilpotent.
For every integer $n\ge1$, the logarithmic probability of the joint
event ``first excess occurs and at least $n$ future visits to
$I_{X,j}$ follow'' is precisely
\begin{equation}\label{eq:fe-operator-survival}
 p_{X,j,n}(L)
 :=\sum_{y\in W_X}\eta_{X,L}(y)
       {\bf1}_{\{N^+_{X,j}(y)\ge n\}}
 =\eta_{X,L}E_{X,j}Q_{X,j}^{\,n-1}{\bf1}.
\end{equation}
Consequently,
\begin{equation}\label{eq:fe-operator-excess}
 \boxed{\quad
 \mathcal E_X(L;K)=\sum_{j\ge0}2^{-j}
      \sum_{n>K_{X,j}}
          \eta_{X,L}E_{X,j}Q_{X,j}^{\,n-1}{\bf1}.
 \quad}
\end{equation}
Here $\mathbf1$ denotes the all-one vector on the appropriate
finite shell. Thus the new open input is an estimate on a
\emph{specific first-excess entrance measure transported by the
actual return operator}, not equidistribution of unused parity bits.
\end{proposition}
\begin{proof}
On a terminating deterministic orbit, at least $n$ future
$j$-shell visits occur if and only if its first future shell entry
exists and is followed by $n-1$ further returns to the same shell.
This is exactly the event counted by the product of kernels on the
right of \eqref{eq:fe-operator-survival}.  Finitely many states occur,
so integration against the entrance law commutes with the kernel
products.  Finally apply the integer tail-sum identity
$(N-K)_+=\sum_{n>K}{\bf1}_{\{N\ge n\}}$ to obtain
\eqref{eq:fe-operator-excess}.
\end{proof}

\begin{corollary}[A polynomial, not geometric, return-tail criterion]
\label{cor:fe-polynomial}
Use the allowances \eqref{eq:fe-allowances}. Suppose there are
constants $A<\infty$ and $\kappa>0$, independent of $X,j,n$, such
that, for all sufficiently large $X$, all $j$ with $I_{X,j}\ne\varnothing$,
and every integer $n>K_{X,j}$,
\begin{equation}\label{eq:fe-polynomial-assumption}
 \sum_{y\in W_X}\eta_{X,L_X}(y)
        {\bf1}_{\{N^+_{X,j}(y)\ge n\}}
 \le A\,p_X(L_X)\left(\frac{K_{X,j}}n\right)^{1+\kappa}.
\end{equation}
Then
\[
 T_X=O\bigl(X(\log X)^{\sigma-\varepsilon_0}\bigr)=o(X).
\]
The polynomial survival bound is required \emph{only after the first
excess crossing}; it is not a claim of renewal independence or of a
uniform second moment on all orbits.
\end{corollary}

\begin{proof}
For a nonnegative integer-valued $N$ and integer $K\ge1$,
$(N-K)_+=\sum_{n>K}{\bf1}_{\{N\ge n\}}$.  Using
\eqref{eq:fe-polynomial-assumption} and a decreasing-function
integral bound gives
\[
 \begin{aligned}
 \sum_{y\in W_X}\eta_{X,L_X}(y)
      \bigl(N^+_{X,j}(y)-K_{X,j}\bigr)_+
 &\le A p_X(L_X)K_{X,j}^{1+\kappa}
          \sum_{n>K_{X,j}}n^{-1-\kappa}\\
 &\le \frac{A}{\kappa}p_X(L_X)K_{X,j}.
 \end{aligned}
\]
Therefore
\[
\mathcal E_X(L_X;K)\le\frac{A}{\kappa}p_X(L_X)
  \sum_j2^{-j}K_{X,j}.
\]
The right-hand side of \eqref{eq:fe-sandwich} is
$O((\log X)^{\sigma-\varepsilon_0})$, proving the claim.
\end{proof}

\begin{remark}[Research status and novelty boundary]
\label{rem:fe-scope}
Theorems~\ref{thm:first-excess-identity} and
\ref{thm:fe-free-budget} are deterministic stopping-line accounting,
combined with the manuscript's stated tightness estimate.  They do
\emph{not} improve the Terras/Tao/GGM mixing estimates, prove the
polynomial return-tail assumption \eqref{eq:fe-polynomial-assumption},
or establish $T_X=o(X)$ unconditionally.  Their value for the
present research programme is the sharper target:
\emph{transport the small first-excess entrance measure through only
very late, scale-weighted return visits}.  This replaces demands for
uniform moment bounds across every return history by a targeted
survival estimate for a rare stopped entrance law.  Whether that
estimate follows from the actual Syracuse dynamics remains open.
\end{remark}
